\documentclass[10pt,a4paper]{amsart}
\usepackage{amsmath,amssymb,amsthm,graphicx}
% Retained minimal preamble for the Stacks Project chapter more-algebra.tex (More on Algebra).
% The chapter file pinned-original-source.tex carries only \input{preamble} on its first line, so no
% preamble text lives inside the chapter itself.  The Stacks Project uses one single shared preamble for
% every chapter; the snapshot stages it identically as sources/stacks-chow/preamble.tex and
% sources/stacks-etale-cohomology/preamble.tex.  The block below is copied mechanically, line for line,
% from sources/stacks-chow/preamble.tex lines 14-18 (the Xy-pic setup) and 154-217 (the native theorem
% declarations together with the mathematical macros).  Nothing was added, removed or redefined; the
% wrapper class and the amsmath/amssymb/amsthm/graphicx base packages are supplied by the assembler.
% Chapter: more-algebra.tex -- More on Algebra.
% For commutative diagrams we use Xy-pic
\usepackage[all]{xy}

% We use 2cell for 2-commutative diagrams.
\xyoption{2cell}

% Theorem environments.
%
\theoremstyle{plain}
\newtheorem{theorem}[subsection]{Theorem}
\newtheorem{proposition}[subsection]{Proposition}
\newtheorem{lemma}[subsection]{Lemma}

\theoremstyle{definition}
\newtheorem{definition}[subsection]{Definition}
\newtheorem{example}[subsection]{Example}
\newtheorem{exercise}[subsection]{Exercise}
\newtheorem{situation}[subsection]{Situation}

\theoremstyle{remark}
\newtheorem{remark}[subsection]{Remark}
\newtheorem{remarks}[subsection]{Remarks}

\numberwithin{equation}{subsection}

% Macros
%
\def\lim{\mathop{\mathrm{lim}}\nolimits}
\def\colim{\mathop{\mathrm{colim}}\nolimits}
\def\Spec{\mathop{\mathrm{Spec}}}
\def\Hom{\mathop{\mathrm{Hom}}\nolimits}
\def\Ext{\mathop{\mathrm{Ext}}\nolimits}
\def\SheafHom{\mathop{\mathcal{H}\!\mathit{om}}\nolimits}
\def\SheafExt{\mathop{\mathcal{E}\!\mathit{xt}}\nolimits}
\def\Sch{\mathit{Sch}}
\def\Mor{\mathop{\mathrm{Mor}}\nolimits}
\def\Ob{\mathop{\mathrm{Ob}}\nolimits}
\def\Sh{\mathop{\mathit{Sh}}\nolimits}
\def\NL{\mathop{N\!L}\nolimits}
\def\CH{\mathop{\mathrm{CH}}\nolimits}
\def\proetale{{pro\text{-}\acute{e}tale}}
\def\etale{{\acute{e}tale}}
\def\QCoh{\mathit{QCoh}}
\def\Ker{\mathop{\mathrm{Ker}}}
\def\Im{\mathop{\mathrm{Im}}}
\def\Coker{\mathop{\mathrm{Coker}}}
\def\Coim{\mathop{\mathrm{Coim}}}

% Boxtimes
%
\DeclareMathSymbol{\boxtimes}{\mathbin}{AMSa}{"02}

%
% Macros for moduli stacks/spaces
%
\def\QCohstack{\mathcal{QC}\!\mathit{oh}}
\def\Cohstack{\mathcal{C}\!\mathit{oh}}
\def\Spacesstack{\mathcal{S}\!\mathit{paces}}
\def\Quotfunctor{\mathrm{Quot}}
\def\Hilbfunctor{\mathrm{Hilb}}
\def\Curvesstack{\mathcal{C}\!\mathit{urves}}
\def\Polarizedstack{\mathcal{P}\!\mathit{olarized}}
\def\Complexesstack{\mathcal{C}\!\mathit{omplexes}}
% \Pic is the operator that assigns to X its picard group, usage \Pic(X)
% \Picardstack_{X/B} denotes the Picard stack of X over B
% \Picardfunctor_{X/B} denotes the Picard functor of X over B
\def\Pic{\mathop{\mathrm{Pic}}\nolimits}
\def\Picardstack{\mathcal{P}\!\mathit{ic}}
\def\Picardfunctor{\mathrm{Pic}}
\def\Deformationcategory{\mathcal{D}\!\mathit{ef}}
\usepackage{cleveref}
\crefname{originalremark}{original remark}{original remarks}
\Crefname{originalremark}{Original remark}{Original remarks}
\makeatletter
\newlabel{remark-construction}{{33882}{source}{}{sourcepointer-1}{}}
\newlabel{remark-construction@cref}{{[originalremark][0][]33882}{[1][0][]source}{}{}{}}
\makeatother
\title{The automorphism group of a normal extension acts transitively on the maximal ideals of the integral closure}
\author{The Stacks Project Authors}
\date{stacks-project:more-algebra.tex:local-snapshot-20261012; GNU FDL 1.2 or later}
\begin{document}
\maketitle
\setcounter{subsection}{1}
\noindent\textit{Editorial note.} All delivered mathematical text is the author's own, copied from the pinned Stacks Project chapter file more-algebra.tex; no statement or proof step was rewritten, reordered or removed.\par\medskip
\noindent\textit{Editorial note.} The chapter file pinned-original-source.tex carries only \textbackslash{}input\{preamble\} on its first line, so no preamble text lives inside the chapter. The retained file sources/stacks-more-algebra/retained-preamble.tex is a mechanical, line-for-line copy of the project's single shared preamble as staged in the snapshot; nothing was added, removed or redefined and no package was added.\par\medskip
\noindent\textit{Editorial note.} The source URL is the source-level Stacks Project URL, as required when more than one card is produced from one source; the per-claim tag is carried by claim\_key.\par\medskip
\noindent\textit{Editorial note.} The single \textbackslash{}ref pointer inside the counted proof span refers to the construction remark and is recorded as a non-load-bearing navigation pointer.\par\medskip
\noindent\textit{Allowed standard prerequisites (editorial).}\begin{itemize}
\item The rings K\_1, L\_1, A\_1 and B\_1 attached to an extension of discrete valuation rings and a normal extension. the construction invoked at statement lines 34612-34614. Fixes the objects on which G acts and among whose maximal ideals transitivity is asserted.
\item The reduced tensor product (L tensor\_K K\_1)\_red and the canonical map from it in characteristic p by a large p-power. used at proof lines 34623-34630. Is the device that lifts the G-invariant element b' back to L tensor\_K K\_1.
\item Normality of the discrete valuation ring B and the identity K = (K\_1)\textasciicircum{}G for a normal extension. used at proof lines 34631-34632. Forces the lifted invariant element to lie in B, which is what produces the contradiction.
\end{itemize}\medskip
\ifdefined\hypertarget\hypertarget{sourcepointer-1}{}\fi
\noindent\textit{Source reference.} Original remark 33882: sources/stacks-more-algebra/pinned-original-source.tex, source lines 33882--33882. Same-chapter navigation pointer; the statement it names is standard for the section and the argument that consumes it stays inside the counted proof span.\par\medskip
\begin{lemma}
\label{lemma-galois-relative}
Let $A \to B$ be an extension of discrete valuation rings with fraction fields
$K \subset L$. Let $K_1/K$ be a normal extension. Say
$G = \text{Aut}(K_1/K)$. Then $G$ acts on the rings $K_1$, $L_1$,
$A_1$ and $B_1$ of Remark \ref{remark-construction}
and acts transitively on the set of maximal ideals of $B_1$.
\end{lemma}

\begin{proof}
Everything is clear apart from the last assertion. If there are two or
more orbits of the action, then we can find an element $b \in B_1$
which vanishes at all the maximal ideals of one orbit and has residue
$1$ at all the maximal ideals in another orbit. Then
$b' = \prod_{\sigma \in G} \sigma(b)$ is a $G$-invariant element of
$B_1 \subset L_1 = (L \otimes_K K_1)_{red}$ which is in some maximal
ideals of $B_1$
but not in all maximal ideals of $B_1$. Lifting it to an element of
$L \otimes_K K_1$ and raising to a high power we obtain a $G$-invariant
element $b''$ of $L \otimes_K K_1$ mapping to $(b')^N$ for some $N > 0$;
in fact, we only need to do this in case the characteristic is $p > 0$ and
in this case raising to a suitably large $p$-power $q$ defines a
canonical map $(L \otimes_K K_1)_{red} \to L \otimes_K K_1$.
Since $K = (K_1)^G$ we conclude that $b'' \in L$. Since $b''$ maps
to an element of $B_1$ we see that $b'' \in B$ (as $B$ is normal).
Then on the one hand it must be true that $b'' \in \mathfrak m_B$
as $b'$ is in some maximal ideal of $B_1$ and on the other hand it
must be true that $b'' \not \in \mathfrak m_B$ as $b'$ is not in
all maximal ideals of $B_1$. This contradiction finishes the proof of the
lemma.
\end{proof}

\end{document}
